\section{Probe-Guided Expert Handoff}
\label{sec:signal}

\subsection{The local--expert loop}
\label{sec:system}
Our system pairs a local full-duplex speech model with a stronger
external expert (Figure~\ref{fig:hero}). The local model processes
incoming speech and decides when to respond. At that transition,
the system chooses whether to continue with the local answer or
request expert help. On the expert path, an audio uplink transcribes
the question, the expert returns a text answer, and the local
talker voices the prepared relay text. The expert request and relay
add cost and waiting time. Our goal is to improve answer accuracy
under a limited expert-call budget by selecting questions that
benefit from expert help. We use local-answer failure risk to guide
this selection; the benefit still depends on whether the expert
repairs the error and the relay preserves its answer.

\subsection{A probe for the handoff decision}
\label{sec:probe}
The local model already computes hidden states as it processes the
question. If these states carry information about its ability to
answer, a small classifier can read that signal without generating
a separate self-evaluation. We use a linear \emph{probe} to test
whether the existing representations contain a directly readable
failure signal.

Offline, we collect local answers to calibration questions and use
a judge to assign correctness labels. Each label is paired with
hidden-state features captured at the chosen read point. We fit a
logistic probe to predict failure from these features while keeping
the speech model frozen. For features $\phi(x)$ from input context
$x$, the fitted weights $w,b$ give the risk score
\begin{equation}
 s(x)=\sigma\!\left(w^\top\phi(x)+b\right).
 \label{eq:gate}
\end{equation}
At inference, the probe reads the current features without a
reference answer or an online judge. Scores above a
language-specific threshold $\tau_{\ell}$ favor escalation;
thresholds are set from calibration scores and held fixed during
evaluation (\S\ref{sec:models}). An auxiliary dialogue-act head
restricts escalation to information requests, filtering turns such
as acknowledgments and stop commands (Appendix~\ref{app:actgate}).
